\section{Related Work}
%%%%%%%%%%%%%%%%%%%%%%%%%%%%%%%%%%%%%%%%%%%%%%%%%%%%%%%%%%%
% Blossom algorithm
% Sequential versions
\textbf{Blossom Algorithms.}
The blossom algorithm~\cite{edmonds1965paths} is foundational for computing maximum matchings in general graphs. Subsequent work improves its sequential performance through optimized data structures and algorithmic refinements~\cite{gabow1976efficient, gabow1984efficient, gabow1990data, gabow1991faster, blum2015maximum}. Despite these advances, Edmonds’ algorithm remains the dominant practical framework.
% Parallel versions
Parallel implementations have also been explored~\cite{Stanford_2016, Wayne_2024, fan2025x}. Among them, X-Blossom removes recursion and reformulates blossom handling, enabling GPU-oriented implementations of maximum matching in general graphs.

%%%%%%%%%%%%%%%%%%%%%%%%%%%%%%%%%%%%%%%%%%%%%%%%%%%%%%%%%%%
% BFS
\noindent
\textbf{BFS.}
As a fundamental graph traversal primitive, BFS has been studied extensively.
Most systems implement BFS in a level-synchronous manner, advancing a frontier level by level~\cite{bfs_levelsync}.
Direction-optimizing BFS (DOBFS) improves performance by switching between top-down push and bottom-up pull to reduce redundant edge examinations~\cite{beamer2011searching,beamer2013direction}. Subsequent work improves frontier representations (sparse vs.\ dense)~\cite{bfs_standford,shun2013ligra,shun2015ligraplus}, reduces contention (e.g., bitmap visited states and fewer atomics)~\cite{bfs_agarwal}, and enhances locality and scalability via frontier reordering and NUMA-aware scheduling~\cite{bfs_locality,bfs_numa,zhang2015numa,arai2016rabbit}. Related implementations appear across distributed systems (Pregel~\cite{malewicz2010pregel}, GraphLab~\cite{low2012graphlab}, PowerGraph~\cite{gonzalez2012powergraph}, Gemini~\cite{zhu2016gemini}), shared-memory CPU frameworks (Galois~\cite{nguyen2013galois}, Ligra/Ligra+~\cite{shun2013ligra,shun2015ligraplus}, Julienne~\cite{ligra_julienne}, Aspen~\cite{ligra_aspen}), and GPU frameworks (Gunrock~\cite{wang2016gunrock}, Groute~\cite{ben2017groute}, Lux~\cite{jia2017lux}, SEP-Graph~\cite{wang2019sep}). Finally, several systems auto-tune these choices for hardware- and graph-specific implementations~\cite{zhang2018graphit,graphit_ordered2020,graphit_gpu,graphit_novelarch}.

%%%%%%%%%%%%%%%%%%%%%%%%%%%%%%%%%%%%%%%%%%%%%%%%%%%%%%%%%%%
% SSSP
\noindent
\textcolor{blue}{\textbf{SSSP and MSSP.}
SSSP extends BFS-style traversal to weighted graphs and is commonly optimized around the tension between priority ordering, redundant work, and exposed parallelism.
Classical approaches include Dijkstra's priority-ordered algorithm and parallel Dijkstra variants~\cite{dijkstra1959note,crauser1998parallelization}, while $\Delta$-stepping uses bucketed relaxation to trade exact priority order for parallel work~\cite{meyer2003delta}.
Subsequent multicore SSSP techniques refine this trade-off through radius/stepping variants, DSMR, relaxed priority schedulers such as SprayList and MultiQueues, and Wasp's priority-aware work stealing~\cite{blelloch2016radius,maleki2016dsmr,dong2021efficient,alistarh2015spraylist,rihani2015multiqueues,postnikova2022multiqueues,dantonio2025wasp}.
SSSP implementations also appear in GAP, GBBS, Galois, Ligra, Gunrock, and GraphIt, where scheduling choices balance locality, synchronization, load balance, and redundant relaxation~\cite{azad2020gap,dhulipala2021gbbs,nguyen2013galois,shun2013ligra,wang2016gunrock,zhang2018graphit}.
On GPUs, work-efficient methods such as Workfront Sweep, Near-Far, and Bucketing reduce redundant relaxation while exposing parallel work~\cite{davidson2014work}; later GPU SSSP work improves scheduler and worklist quality through asynchronous dynamic $\Delta$-stepping, bucket-aware execution, dynamic-stepping heuristics, and multi-level multi-queue designs~\cite{wang2021fastsssp,zhang2023bucketaware,yu2025heuristic,hu2026mlmq}.}

%%%%%%%%%%%%%%%%%%%%%%%%%%%%%%%%%%%%%%%%%%%%%%%%%%%%%%%%%%%
% MSSP
\textcolor{blue}{MSSP has been studied as both a planar-graph primitive and batched SSSP in general graph analytics.
Planar-graph MSSP techniques compactly represent shortest-path trees for all boundary-face sources and compute distances from those roots ~\cite{klein2005multiple}.
This line of work has also been extended to embedded-graph MSSP and linear-time algorithms for unit-weight planar MSSP~\cite{cabello2013multiple,eisenstat2013linear}.
Beyond these specialized techniques, general graph analytics often treats MSSP as batched shortest-distance computation, using batched Dijkstra and Bellman--Ford variants~\cite{then2017batched}.
This batched view also aligns with frameworks supporting MSSP through frontier-based traversal operators and scheduling abstractions~\cite{shun2013ligra,wang2016gunrock,zhang2018graphit}.}

%%%%%%%%%%%%%%%%%%%%%%%%%%%%%%%%%%%%%%%%%%%%%%%%%%%%%%%%%%%
% CPU and GPU analysis => Debunking
\noindent
\textbf{The Platform Selection.}
%The CPU versus GPU detbate has remained heated for more than two decades. 
% GPU is fast
With programmable GPUs, many domain experts rewrote performance-critical kernels for massive parallelism and reported striking speedups across diverse domains such as networking ~\cite{gpu_network_1,gpu_network_2}, databases analytics ~\cite{gpu_db_1,gpu_db_2}, and classic algorithmic primitives ~\cite{gpu_algo_1,gpu_algo_2}. 
However, early claims of substantial GPU performance advantages ~\cite{gpuIsFast_che,gpuIsFast_ryoo} motivated systematic CPU-GPU comparisons. 
% CPU
Meanwhile, CPUs have continued to evolve, adopting wider SIMD vectors~\cite{cpu_simd_1,cpu_simd_2} and integrating higher-bandwidth memory technologies~\cite{cpu_hbm_1,cpu_hbm_2}. 
Recent CPU-centric studies also report strong performance and energy efficiency across modern scientific and memory-bound workloads~\cite{cpu_perf_1}, and in some cases, CPU kernels can even rival or outperform GPU counterparts~\cite{cpu_over_gpu_1,cpu_over_gpu_2,cpu_over_gpu_3}.
% Platform Selecition
%These ongoing comparisons have motivated more rigorous CPU–GPU evaluations that emphasize architecture-aware optimization and insightful interpretation metrics, rather than relying on raw runtime speedups alone.
A study has reexamined early reports of dramatic GPU gains~\cite{gpuIsFast_che,gpuIsFast_ryoo}, showing that many such speedups shrink markedly when CPU and GPU implementations are optimized to a comparable level and evaluated fairly with respect to architectural constraints~\cite{lee2010debunking}.
The importance of including host--device data transfer overheads in end-to-end measurements is also emphasized ~\cite{gregg2011transfer}. 
Power and energy efficiency have also been studied to inform system design~\cite{gpu_energy}, while newer interpretation metrics (e.g., peak-based crossover) help determine when GPUs outperform CPUs for batch workloads~\cite{peek_metrics2025}.
Finally, standardized benchmark suites support comparable cross-platform evaluation, including GAP for graph analytics~\cite{benchmark_gap,benchmark_gap_eval} and MLPerf for machine learning~\cite{benchmark_mlperf,benchmark_mlperfhpc}. 

%These related efforts focus on software-level optimizations and are largely case-specific. In contrast, our profiling metrics capture critical architectural factors governing program execution, quantify hardware resource utilization, and assess workload suitability for either multicore CPUs or GPUs.
All of the aforementioned related efforts focus on software-level optimizations and are largely case-specific. In contrast, our profiling metrics capture the critical architectural factors governing program execution, quantifying hardware resource utilization levels and enabling a systematic assessment of workload suitability for either multicore CPUs or GPUs.
